% !TeX root = ../main.tex
% 四章架構的通用研究骨架；背景與方法須依作者實際研究填寫。
\chapter{Introduction}
\label{ch:introduction}
\thesisplaceholder{Introduce the biomedical problem, relevant evidence, research questions, and the organization of this thesis.}

\section{Biomedical Problem and Motivation}
% 說明具體臨床或生物學問題、現有流程、目標族群及重要性。
% 指定使用者、決策時點與資訊可得性；用核實來源支持需求，不以模型名稱代替問題。

\section{Background and Related Work}
% 介紹理解研究所需的背景，再按研究問題、資料、方法、評估與限制整理原始文獻。
% 比較現行實務及簡單而強的基準；不直接比較任務、族群或評估設定不相容的分數。
% 不預設資料模態或生物機制；僅加入實際研究相關、已閱讀核實的內容。

\section{Research Gap, Questions, and Objectives}
% 說明文獻仍未解決的問題；區分未研究、證據不足與已得到負結果。
% 定義輸入、輸出、分析單位、主要終點與成功條件；分開確認性假設及探索性問題。

\section{Scope and Contributions}
% 界定本研究處理及不處理的情境；每項實際貢獻都應能對應後續章節證據。
% 預計執行、已實作及已驗證須明確區分，不把計畫寫成已有成果。

\section{Thesis Structure}
% 依最終內容概述第二章 Materials and Methods、第三章 Results and Discussion、
% 第四章 Conclusion 如何回答研究問題；引用附錄時同步核對其編號與內容。
